% ============================================================
% study_guide PDF 조판 헤더 (pandoc -H 로 주입)
% ============================================================

% ----- 줄 간격 / 문단 -----
\usepackage{setspace}
\setstretch{1.25}
\setlength{\parskip}{0.55em plus 0.1em}
\setlength{\parindent}{0pt}

% ----- 단어 줄바꿈(영문 긴 식별자) 허용, 과한 hyphen 억제 -----
\usepackage{microtype}
\tolerance=2000
\emergencystretch=2em
\hyphenpenalty=2000
\sloppy

% ----- 리스트 간격/들여쓰기 -----
\usepackage{enumitem}
\setlist{itemsep=2pt, parsep=1pt, topsep=3pt, leftmargin=1.5em}

% ----- 표 행 간격 -----
\renewcommand{\arraystretch}{1.25}

% ----- 섹션 제목 위아래 간격 -----
\usepackage{titlesec}
\titlespacing*{\section}{0pt}{1.4em}{0.6em}
\titlespacing*{\subsection}{0pt}{1.0em}{0.4em}

% ----- 이미지: 최대 크기 제한 + 중앙 (원본보다 확대하지 않음 → 항상 선명, 크기 일관) -----
\usepackage{graphicx}
\usepackage[export]{adjustbox}
\let\origincludegraphics\includegraphics
\renewcommand{\includegraphics}[2][]{%
  \origincludegraphics[#1,max width=0.74\linewidth,max totalheight=0.6\textheight,center]{#2}}

% ----- 코드 / 블록 다이어그램 / 트리 = 인용구식 박스 -----
\usepackage{fancyvrb}
\usepackage{tcolorbox}
\tcbuselibrary{breakable,skins}
\definecolor{codebg}{RGB}{247,248,250}
\definecolor{codeframe}{RGB}{203,209,217}

% (1) 언어 지정 코드블록 = pandoc Shaded 환경 (highlight 미사용 대비 가드)
\makeatletter
\@ifundefined{Shaded}{\newenvironment{Shaded}{}{}}{}
\makeatother
\renewenvironment{Shaded}{%
  \begin{tcolorbox}[breakable, enhanced, arc=2pt, boxrule=0.4pt,
    colback=codebg, colframe=codeframe,
    left=7pt, right=7pt, top=4pt, bottom=4pt,
    before skip=7pt, after skip=9pt]%
  \footnotesize}{%
  \end{tcolorbox}}

% (2) plain 코드블록·ASCII 다이어그램·트리 = pandoc verbatim 환경 → 박스
%     fancyvrb Verbatim 을 tcolorbox 로 감싸 재정의
\renewenvironment{verbatim}{%
  \VerbatimEnvironment
  \begin{tcolorbox}[breakable, enhanced, arc=2pt, boxrule=0.4pt,
    colback=codebg, colframe=codeframe,
    left=7pt, right=7pt, top=3pt, bottom=3pt,
    before skip=7pt, after skip=9pt]%
  \begin{Verbatim}[fontsize=\footnotesize]}{%
  \end{Verbatim}\end{tcolorbox}}

\fvset{baselinestretch=1.0}
